% !TeX spellcheck = ru_RU
% !TEX root = vkr.tex
\section{Обзор и выводы}
\label{sec:relatedworks}

PySATL-expert формирует рейтинг распределений по характеристикам выборки
и статистикам критериев. Работа NIST~\cite{nistDistributionNN} показывает
другой вариант распознавания: нейронная сеть получает представление
ядерной оценки плотности. Для нашей работы это исследовательский аналог,
но его результаты не доказывают преимущество над существующей моделью Expert.

Программные аналоги fitdistrplus~\cite{fitdistrplusOverview}
и Phitter~\cite{phitter2025} объединяют оценивание параметров,
диагностику подгонки и сравнение кандидатов. Поэтому результат Expert
должен содержать не только рейтинг, но и параметры, результаты проверок
и ограничения анализа: окончательный выбор остаётся за статистиком.

Документация SciPy~\cite{scipyGof} описывает параметрическую bootstrap-проверку:
неизвестные параметры повторно оцениваются на каждой модельной выборке.
Следовательно, калибровку необходимо связывать со способом оценивания;
таблицу для заранее известных параметров нельзя автоматически применять
после их подгонки. Повторное использование готового банка требует
отдельного обоснования для соответствующей процедуры.

Обзор тестов экспоненциальности~\cite{exponentialityReview} показывает,
что мощность специализированных критериев зависит от альтернативы.
Наличие большого каталога проверок само по себе не гарантирует
качества рекомендаций.

\subsection*{Выводы}
\label{subsec:review-conclusions}
Из обзора следуют три решения: отделить настраиваемый анализ от входов
обученной модели; сохранять способ оценивания и условия калибровки;
сравнивать варианты рекомендаций по качеству первого места, короткого
списка и времени работы. Сначала будут сопоставлены представления
признаков при одинаковых модели и разбиении данных, затем --- модели.
Это обосновывает три задачи ВКР. Полная подборка источников и подробные
заметки сохранены в \texttt{docs/MATERIALS.md}.
